\section{Length of Names and Use of Abbreviations}\label{sec:LengthOfNamesAndUseOfAbbreviations}
The \bdq{The \Java Language Specification} does not impose a limit for names of program elements: \bdq{An \textit{identifier} is an unlimited-length sequence of \textit{Java letters} and \textit{Java digits}, the first of which must be a \textit{Java letter}.} \autocite{ORACLE_DOC_LANGUAGE_SPECIFICATION:Identifiers}

But of course Java programs do have a maximum length for identifiers (because the compiled code may not exceed a given length), although this is higher than any practically usuable value. I think that a name for a module, package, class, method, field, constant, argument, or local variable that is longer than 100 or perhaps 150~characters is definitely too long to be useful.

In general, you should avoid abbreviations and acronyms as identifiers for program elements if those are not \textit{very well known} – refer to \tqfullref{sec:GeneralAccords} where I already discussed the need for meaningful names.

The length of a name has very little to no relevant influence on the size and speed of the compiled code. And with the code completion features provided by modern IDEs, the programmer may not even have to type them when referencing the named elements~…

This means that you may never sacrifice the meaningfulness of a long name for typing speed – never ever!

If you think you will be more productive when using short names (because you can type them faster~…), you are wrong! When you have the feeling that you have to write code faster, learn typewriting\footnote{I discussed that already on page \pageref{ref:Typewriting}.}. It will also help you to write all the other stuff you have to deliver in addition to your code (documentation, meeting notes, emails, specification documents,~…).

%==============================================================================
% End of file!!
%==============================================================================
